\protect \section *{Розділ~\ref *{sec:particle-motion}. \nameref *{sec:particle-motion}}
\begin{Solution}{3.{1}}
		$r_{Le}\approx3{,}4\times10^{-3}$~см $<\lambda_{De}\approx7{,}4\times10^{-3}$~см.
	
\end{Solution}
\begin{Solution}{3.{2}}
		$v_E=cE/B=10^{5}$~см/с (1~км/с), однакова для електрона й~протона.
	
\end{Solution}
\begin{Solution}{3.{3}}
		електрон релятивістський ($\gamma\approx3$, $v\approx0{,}94c$); $r_L\approx5$~км,
	$2\pi/\omega_c\approx2\times10^{-4}$~с, $\tau_b\approx0{,}4$~с, $\tau_d\approx15$~хв.
	
\end{Solution}
\begin{Solution}{3.{5}}
		$v_\perp$ стала, а~вихід настає за $v_\parallel/v_\perp=\cot\theta_c=2$, тобто коли $L$ зменшиться вдвічі:
	$t\approx2{,}5\times10^8$~с ($\approx8$~років), $\mathcal E\approx2{,}5$~кеВ.
	
\end{Solution}
\begin{Solution}{3.{6}}
		$\sin\theta_c=\sqrt{B_{\min}/B_{\max}}=1/2$, $\theta_c=30^\circ$; частка $1-\cos\theta_c\approx13\,\%$.
	
\end{Solution}
\begin{Solution}{3.{7}}
		$E=v_EB/c\approx33$~СГС $\approx10^4$~В/см; від сорту іонів не залежить.
	
\end{Solution}
\begin{Solution}{3.{8}}
		$\omega=2c\varphi_0/(Ba^2)\approx2\times10^6$~рад/с.
	
\end{Solution}
\begin{Solution}{3.{9}}
		$v_d\approx2$~км/с, $t\approx1{,}5\times10^{-4}$~с.
	
\end{Solution}
\begin{Solution}{3.{10}}
		$\omega_{ce}=2\pi\cdot126$~ГГц відповідає $B\approx4{,}5$~Тл, $R=R_0\cdot5/4{,}5\approx1{,}83$~м; точка зміщується до $R\approx2{,}0$~м.
	
\end{Solution}
\begin{Solution}{3.{11}}
		$T_i=20$~еВ.
	
\end{Solution}
